\setcounter{section}{4}
\section{Bauelemente}
\subsection{Kondensator I}

\begin{sol}{EA101}
  Die Einheit Farad (F) musst Du Dir merken.
\end{sol}

\begin{sol}{EC201}
  Legen wie eine Gleichspannung an einen Kondensator, so fließt der Strom zunächst in den Kondensator, deshalb ist die Spannung zunächst 0. Relativ schnell wir der Kondensator allerdings geladen sein, und dann steigt die Spannung schnell an. Dies entspricht dem Kurvenverlauf der Musterantwort (A).
\end{sol}


\begin{sol}{EC202}
  Du kannst es Dir vielleicht folgendermaßen merken:
  Bei Frequenz \SI{0}{\hertz} (Gleichstrom) leitet ein Kondenstor überhaupt nicht. Erst mit steigender Frequenz wird er überhaupt leitend.
\end{sol}


\begin{sol}{EC203}
 Es ist nach gefragt, was die Kapazität \qq{verringert}. Die Kapazität hängt im wesentlichen von Plattenabstand, Plattenfläche und Dielektrikum ab. Nur bei Antwort (A) stimmt die Richtung. 
\end{sol}


\begin{sol}{EC204}
  Siehe Frage \qref{EC203}.
\end{sol}

\begin{sol}{EC205}
  Siehe Frage \qref{EC203}.
\end{sol}

\begin{ohmchapter}
\end{ohmchapter}

\subsection{Spule I}

\begin{ohmchapter}
\end{ohmchapter}

\subsection{Übertrager I}

\begin{ohmchapter}
\end{ohmchapter}

\subsection{Diode I}

\begin{ohmchapter}
\end{ohmchapter}

\subsection{Transistor I}

\begin{ohmchapter}
\end{ohmchapter}